\chapter{Anleitung: Inbetriebnahme}
\label{ch:anleitung}

Dieses Kapitel beschreibt die grundlegenden Schritte zur Inbetriebnahme von
WikiPod. Die Anleitung ist insbesondere als Einstiegspunkt für Folgeprojekte
gedacht. Weiterführende Informationen zur Projektstruktur und zu einzelnen
Komponenten befinden sich in der README-Datei des Repositorys.

\section{Voraussetzungen}

Für die lokale Ausführung wird eine Python-Umgebung sowie Docker mit Docker
Compose benötigt. Die Anwendung verwendet OpenSearch für die Speicherung und
Abfrage der erzeugten Vektorrepräsentationen. Für die Dokumentengrundlage wird
außerdem eine lokale Wikipedia-ZIM-Datei benötigt. Diese ist aufgrund ihrer
Größe nicht Bestandteil des Repositorys und muss separat bereitgestellt werden.

Für den produktionsnahen Betrieb wurde im Projekt ein Raspberry Pi~5 mit
16\,GB Arbeitsspeicher und einer 1-TB-SSD verwendet. Die grundlegenden
Entwicklungs- und Testschritte können jedoch auch auf einem gewöhnlichen
Entwicklungsrechner durchgeführt werden.

\section{Installation}

Nach dem Klonen des Repositorys wird zunächst eine virtuelle Python-Umgebung
erstellt und anschließend das Projekt einschließlich der
Entwicklungsabhängigkeiten installiert.

\begin{lstlisting}[language=bash]
git clone https://github.com/niklasbelieres/wikipod-rag
cd wikipod-rag

python -m venv .venv
source .venv/bin/activate
pip install -e ".[dev]"
\end{lstlisting}

Anschließend muss eine Wikipedia-ZIM-Datei lokal bereitgestellt werden. Der
Pfad zu dieser Datei wird über den Eintrag \texttt{paths.zim\_file} in der
jeweiligen Konfigurationsdatei festgelegt.

\section{Lokales Setup (Entwicklung)}

OpenSearch wird lokal über Docker Compose gestartet:

\begin{lstlisting}[language=bash]
docker compose up -d
\end{lstlisting}

Ob OpenSearch erreichbar ist, kann anschließend beispielsweise mit folgendem
Aufruf geprüft werden:

\begin{lstlisting}[language=bash]
curl http://localhost:9200
\end{lstlisting}

Wenn die konfigurierte ZIM-Datei vorhanden und OpenSearch gestartet ist, kann
die Indexierung im Entwicklungsmodus ausgeführt werden:

\begin{lstlisting}[language=bash]
WIKIPOD_ENV=dev python -m wikipod.cli index
\end{lstlisting}

Nach erfolgreicher Indexierung kann eine Anfrage an das System gestellt werden:

\begin{lstlisting}[language=bash]
WIKIPOD_ENV=dev python -m wikipod.cli query \
  "What is the Catholic Church?"
\end{lstlisting}

Soll ausschließlich das Retrieval untersucht werden, ohne eine Antwort durch
das lokale Sprachmodell zu erzeugen, kann die Option
\texttt{--chunks-only} verwendet werden:

\begin{lstlisting}[language=bash]
WIKIPOD_ENV=dev python -m wikipod.cli query \
  --chunks-only "What is the Catholic Church?"
\end{lstlisting}

\section{Konfiguration}

Die Konfiguration des Systems befindet sich im Verzeichnis
\texttt{config/}. Unterschiedliche Konfigurationen werden für Entwicklung,
Evaluation, Tests auf dem Raspberry Pi und den produktionsnahen Betrieb
verwendet. Über die Umgebungsvariable \texttt{WIKIPOD\_ENV} kann die jeweils
benötigte Laufzeitkonfiguration ausgewählt werden. Wird keine Umgebung
angegeben, verwendet WikiPod standardmäßig die Entwicklungsumgebung
\texttt{dev}.

Insbesondere muss der Pfad zur verwendeten Wikipedia-ZIM-Datei über
\texttt{paths.zim\_file} an die lokale Umgebung angepasst werden. Dadurch
können unterschiedliche Datensätze und Laufzeitumgebungen verwendet werden,
ohne die allgemeine Projektkonfiguration für jeden Versuch verändern zu
müssen.

\section{Deployment auf dem Raspberry Pi}

Für den Betrieb auf dem Raspberry Pi wird das Repository zunächst auf das
Gerät übertragen und analog zur lokalen Installation eingerichtet. Die
Wikipedia-ZIM-Datei muss auf dem lokalen Speichermedium des Raspberry Pi
vorliegen und der entsprechende Pfad in der dafür vorgesehenen
Konfiguration eingetragen werden.

OpenSearch wird auch auf dem Raspberry Pi über Docker Compose betrieben.
Nach dem Start von OpenSearch kann die Indexierung mit der für den jeweiligen
Versuch vorgesehenen Konfiguration ausgeführt werden. Für Pi-spezifische
Testläufe steht beispielsweise die Konfiguration
\texttt{config/pi-test.yaml} zur Verfügung. Dadurch können Experimente auf dem
Raspberry Pi durchgeführt werden, ohne die Standardkonfiguration für die
Entwicklung zu verändern.

\par\smallskip
\noindent\begin{minipage}[t]{\linewidth}
Bei längeren Indexierungsläufen kann zusätzlich das im Projekt enthaltene
Monitoring gestartet werden:

\begin{lstlisting}[language=bash]
python -m wikipod.scripts.monitor_run \
  --out logs/pi_top100_index.csv \
  --interval 10
\end{lstlisting}
\end{minipage}\par\smallskip

Das Monitoring zeichnet unter anderem CPU-, Speicher-, Swap- und
Temperaturwerte sowie den Speicherverbrauch von WikiPod und OpenSearch auf.
Nach Abschluss des Laufs kann das Monitoring mit \texttt{Ctrl+C} beendet und
die erzeugte CSV-Datei ausgewertet werden:

\begin{lstlisting}[language=bash]
python -m wikipod.scripts.plot_metrics \
  logs/pi_top100_index.csv
\end{lstlisting}

\section{Evaluation ausführen}

Für die Retrieval-Evaluation steht das Kommando \texttt{evaluate} zur
Verfügung. Als Eingabe wird ein YAML- oder JSON-Datensatz mit
natürlichsprachlichen Anfragen und den jeweils relevanten Wikipedia-Artikeln
verwendet.

Für die im Bericht eingesetzte Evaluationskonfiguration kann beispielsweise
folgender Aufruf ausgeführt werden:

\begin{lstlisting}[language=bash]
WIKIPOD_ENV=report_eval python -m wikipod.cli evaluate \
  --dataset test/data/eval_queries.yaml
\end{lstlisting}

Die Anzahl der betrachteten Retrieval-Ergebnisse kann bei Bedarf mit
\texttt{--top-k} überschrieben werden. Ergebnisse lassen sich außerdem als
JSON oder CSV beziehungsweise gesammelt in einem Ausgabeverzeichnis speichern:

\begin{lstlisting}[language=bash]
WIKIPOD_ENV=report_eval python -m wikipod.cli evaluate \
  --dataset test/data/eval_queries.yaml \
  --top-k 5 \
  --output-dir logs/report_eval_k5
\end{lstlisting}

Optional können mit \texttt{--use-query-analyzer} die normalisierten Anfragen
des Query Analyzers verwendet und mit \texttt{--use-llm-judge} zusätzliche
LLM-basierte Bewertungen aktiviert werden.

Vor einem neuen Lauf können die automatisierten Tests und die statische
Codeprüfung ausgeführt werden:

\begin{lstlisting}[language=bash]
pytest
ruff check .
\end{lstlisting}

Diese Prüfungen werden zusätzlich bei Pushes und Pull Requests auf den
\texttt{main}-Branch über GitHub Actions ausgeführt.